%% ============================================================
%%  appendix/B_interview_experiences.tex
%%  Appendix B: Student Interview Experiences
%%
%%  CONTRIBUTOR GUIDE:
%%  To add a new experience, copy the \begin{experience}...\end{experience}
%%  block below and fill in the blanks.
%%  Use \input{} to include individual experience files from this appendix.
%%
%%  Template file: see template_interview.tex (in the appendix/ folder)
%% ============================================================

\chapter{Student Interview Experiences}
\label{app:experiences}

\begin{disclaimer}
  All interview experiences below are shared voluntarily by \PrepFusion students.
  Question sets, panel compositions, and processes \textbf{vary each year} and
  may differ from what is documented here. Use these accounts for orientation
  only---always verify current details with the respective institute.
\end{disclaimer}

%% ─────────────────────────────────────────────────────────
\section{How to Contribute Your Experience}
\label{sec:contribute}

If you have recently appeared for an IIT or IISc interview, please contribute
your experience to help future students! Share via the PrepFusion Google Doc:

\begin{center}
  \url{https://docs.google.com/document/d/1WaXfhvyS9vxOzVhDEFRc2AwDmobg-vkDxYCtel89IqU}
\end{center}

Our interns will format your contribution using the \texttt{template\_interview.tex}
file and include it in the next update of this guide.

%% ─────────────────────────────────────────────────────────
\section{Experience Template Reference}
\label{sec:template-ref}

\begin{experience}{[Student Name or Anonymous]}{[Institute] -- [Branch/Program]}

  \textbf{Date:} [Month, Year]
  \hfill\textbf{Duration:} [Duration of interview]

  \vspace{4pt}
  \textbf{Written Test:}
  [Describe the written test format, difficulty, and topics covered.]

  \vspace{4pt}
  \textbf{Interview Panel:} [Professor names / Anonymous]

  \vspace{4pt}
  \textbf{Questions Asked:}
  \begin{pfChecklist}
    \item [Question 1]
    \item [Question 2]
    \item [Question 3 -- add or remove as needed]
  \end{pfChecklist}

  \vspace{4pt}
  \textbf{Key Takeaways / Tips for Future Candidates:}
  \begin{pfActionList}
    \item [Tip 1]
    \item [Tip 2]
  \end{pfActionList}

\end{experience}

%% ─────────────────────────────────────────────────────────
%%  ADD EXPERIENCE FILES BELOW USING \input{}
%%  Example:
%%  \input{appendix/experiences/iisc_mvlsi_2025.tex}
%%  \input{appendix/experiences/iitb_ics_2025.tex}
%% ─────────────────────────────────────────────────────────

%% Placeholder experience entries follow.
%% Replace with real \input{} calls as contributions arrive.

\section{IISc Bangalore}
\label{sec:exp-iisc}

\begin{experience}{Contributed via PrepFusion Community}{IISc Bangalore -- MVLSI}

  \textbf{Written Test:} Standard analog + digital MCQ format similar to GATE
  difficulty. Some questions on control systems.

  \vspace{4pt}
  \textbf{Analog Interview (Prof.\ Aniruddhan):}
  \begin{pfChecklist}
    \item Wien Bridge oscillator --- draw, explain, what sets frequency?
    \item Op-amp with negative feedback --- draw, derive closed-loop gain.
    \item Dominant pole concept in control systems.
    \item MOSFET differential amplifier --- small-signal model and CMRR.
    \item Written test correction on whiteboard.
  \end{pfChecklist}

  \vspace{4pt}
  \textbf{Digital Interview (Prof.\ Janakiraman \& Prof.\ V.\ Vasudevan):}
  \begin{pfChecklist}
    \item Design a 4-bit binary multiplier using basic gates.
    \item Explain Wallace Tree multiplier.
    \item Domain interest check: Are you genuinely interested in digital or analog?
  \end{pfChecklist}

  \textbf{Key Tip:} IISc emphasizes \textbf{fundamentals} over rote formulas.
  Expect follow-up questions on anything you answered incorrectly in the written
  test. Whiteboard confidence is essential.
\end{experience}

%% ─────────────────────────────────────────────────────────
\section{IIT Delhi}
\label{sec:exp-iitd}

%% Placeholder -- replace with \input{} when experiences are received
\begin{experience}{[Placeholder]}{IIT Delhi -- IEC}
  \textit{This section will be filled with community-sourced interview
  experiences. If you appeared for an IIT Delhi IEC interview, please share
  your experience at the link above.}
\end{experience}

%% ─────────────────────────────────────────────────────────
\section{IIT Bombay}
\label{sec:exp-iitb}

\begin{experience}{[Placeholder]}{IIT Bombay -- ICS / ESE}
  \textit{This section will be filled with community-sourced interview
  experiences. If you appeared for an IIT Bombay MS or M.Tech Research
  interview, please share your experience.}
\end{experience}
